% Preamble for "Kannada, Step by Step, Root by Root" (HL00 book format).
% XeLaTeX only. Uses the vendored static Noto Sans Kannada (see _fonts/),
% loaded by relative path so local and CI builds are identical.
%
% Note: the Kannada script is rendered via a \newfontfamily with Script=Kannada
% (XeTeX does the complex-script shaping from the font's OpenType tables). We do
% NOT load a polyglossia Kannada language module — only english as the main
% language — so the build never depends on a gloss-kannada.ldf being present.

\usepackage{fontspec}
\setmainfont{Latin Modern Roman}
\setsansfont{Latin Modern Sans}
\setmonofont{Latin Modern Mono}

% A few transliteration letters used in the cognate boxes (Tamil alveolar ṉ, ṟ)
% are absent from Latin Modern; render them with the bar-under accent.
\usepackage{newunicodechar}
\newunicodechar{ṉ}{\b{n}}
\newunicodechar{ṟ}{\b{r}}
\newunicodechar{ṁ}{\.{m}}
\newunicodechar{ḻ}{\b{l}}
\newunicodechar{ḱ}{\'{k}}
\newunicodechar{₃}{\textsubscript{3}}

\usepackage[table]{xcolor}
\usepackage{tcolorbox}
\tcbuselibrary{skins,breakable}
\usepackage{booktabs}
\usepackage{longtable}
\usepackage{tabularx}
\usepackage{array}
\usepackage[hidelinks]{hyperref}
\usepackage[raggedright]{titlesec}
% Shared deterministic figure styling for every Human Languages book (HL06).
% SVG is the committed source; the book build converts it to same-name PDF
% before XeLaTeX, which keeps TeX Live's dependency closure lean and reproducible.
\usepackage{graphicx}
\graphicspath{{figures/}}

% A block figure fits a box, not only a width. Width alone was enough while
% every figure was a wide etymology route; a stroke-order filmstrip of a tall,
% narrow letter (alef is about four times taller than wide) scaled to the line
% width ran off the page (HL-C443). With both bounds, keepaspectratio scales the
% figure to whichever limit it meets first, so a wide figure is unchanged.
\newcommand{\hlblockfigure}[2]{%
  \par\medskip
  \begingroup
    \centering
    \includegraphics[width=0.92\linewidth,height=0.45\textheight,keepaspectratio]{#1}\par
    \small\itshape #2\par
  \endgroup
  \medskip
}

\newcommand{\hlinlinefigure}[1]{%
  \raisebox{-0.25ex}{\includegraphics[height=1.4em,keepaspectratio]{#1}}%
}


\usepackage{polyglossia}
\setmainlanguage{english}
\setotherlanguage{arabic}

% Vendored Kannada font (../../_fonts/ from <lang>/book/); Script=Kannada drives
% the shaping. \kn{...} = an inline Kannada snippet within English text.
\newfontfamily\kannadafont[
  Path=../../_fonts/,
  Script=Kannada,
  BoldFont=NotoSansKannada-Static.ttf,
  ItalicFont=NotoSansKannada-Static.ttf,
  BoldItalicFont=NotoSansKannada-Static.ttf
]{NotoSansKannada-Static.ttf}
\newcommand{\kn}[1]{{\kannadafont #1}}
\newfontfamily\tamilfont[
  Path=../../_fonts/,
  Script=Tamil,
  BoldFont=NotoSansTamil-Static.ttf,
  ItalicFont=NotoSansTamil-Static.ttf,
  BoldItalicFont=NotoSansTamil-Static.ttf
]{NotoSansTamil-Static.ttf}
\newcommand{\ta}[1]{{\tamilfont #1}}
\newfontfamily\telugufont[
  Path=../../_fonts/,
  Script=Telugu,
  BoldFont=NotoSansTelugu-Static.ttf,
  ItalicFont=NotoSansTelugu-Static.ttf,
  BoldItalicFont=NotoSansTelugu-Static.ttf
]{NotoSansTelugu-Static.ttf}
\newcommand{\te}[1]{{\telugufont #1}}
\newfontfamily\malayalamfont[
  Path=../../_fonts/,
  Script=Malayalam,
  BoldFont=NotoSansMalayalam-Static.ttf,
  ItalicFont=NotoSansMalayalam-Static.ttf,
  BoldItalicFont=NotoSansMalayalam-Static.ttf
]{NotoSansMalayalam-Static.ttf}
\newcommand{\ml}[1]{{\malayalamfont #1}}
\newfontfamily\devanagarifont[
  Path=../../_fonts/,
  Script=Devanagari,
  BoldFont=NotoSansDevanagari-Static.ttf,
  ItalicFont=NotoSansDevanagari-Static.ttf,
  BoldItalicFont=NotoSansDevanagari-Static.ttf
]{NotoSansDevanagari-Static.ttf}
\newcommand{\dv}[1]{{\devanagarifont #1}}
\newfontfamily\arabicfont[
  Path=../../_fonts/,
  Script=Arabic,
  BoldFont=NotoNaskhArabic-Static.ttf,
  ItalicFont=NotoNaskhArabic-Static.ttf,
  BoldItalicFont=NotoNaskhArabic-Static.ttf
]{NotoNaskhArabic-Static.ttf}
\newcommand{\ar}[1]{\textarabic{#1}}

\hypersetup{
  pdftitle={Kannada, Step by Step, Root by Root},
  pdfauthor={Adhithya Rajasekaran},
  pdfsubject={Etymology-driven Kannada curriculum},
  pdfkeywords={Kannada, Dravidian, language learning, etymology}
}
\pdfstringdefDisableCommands{%
  \def\kn#1{#1}% Keep script text while dropping font-only presentation.
  \def\ta#1{#1}%
  \def\te#1{#1}%
  \def\ml#1{#1}%
  \def\dv#1{#1}%
  \def\ar#1{#1}%
  \def\kannadafont{}%
  \def\tamilfont{}%
  \def\telugufont{}%
  \def\malayalamfont{}%
  \def\devanagarifont{}%
  \def\arabicfont{}%
}

% Micro-lessons intentionally end at natural pause points instead of stretching
% their callout boxes to force every page to the same depth.
\raggedbottom
\emergencystretch=2em

\newtcolorbox{cousinweb}{
  breakable, skin=enhanced, colback=yellow!8, colframe=yellow!45!black,
  boxrule=0.5pt, arc=1mm, left=6pt, right=6pt, top=4pt, bottom=4pt,
  fonttitle=\bfseries, title={The word, taken apart --- its roots and family}
}
\newtcolorbox{culture}{
  breakable, skin=enhanced, colback=red!5, colframe=red!40!black,
  boxrule=0.5pt, arc=1mm, left=6pt, right=6pt, top=4pt, bottom=4pt,
  fonttitle=\bfseries, title={Why it's said this way}
}
\newtcolorbox{grammarlens}[1][]{
  breakable, skin=enhanced, colback=blue!5, colframe=blue!35!black,
  boxrule=0.5pt, arc=1mm, left=6pt, right=6pt, top=4pt, bottom=4pt,
  fonttitle=\bfseries, title={Grammar Lens}, #1
}
\newtcolorbox{sounds}{
  breakable, skin=enhanced, colback=green!6, colframe=green!30!black,
  boxrule=0.5pt, arc=1mm, left=6pt, right=6pt, top=4pt, bottom=4pt,
  fonttitle=\bfseries, title={The letters in this word}
}
\newtcolorbox{cognates}{
  breakable, skin=enhanced, colback=violet!6, colframe=violet!45!black,
  boxrule=0.5pt, arc=1mm, left=6pt, right=6pt, top=4pt, bottom=4pt,
  fonttitle=\bfseries, title={Across the family --- the same idea, five ways}
}

\titleformat{\chapter}[display]
  {\normalfont\Large\bfseries}{Chapter \thechapter}{10pt}{\huge\raggedright}

% Chapter opening (HL09 §8): what the reader will be able to do by the end.
% Generated from the HL05 capability ledger, never hand-written.
\newenvironment{chapteropening}%
  {\begin{quote}\itshape\small\raggedright}%
  {\end{quote}\par\medskip}

% ---------------------------------------------------------------------------
% Table-of-contents number width (HL-C109).
%
% `book.cls` reserves 1.5em for a chapter number in the table of contents:
%
%     \newcommand*\l@chapter[2]{... \setlength\@tempdima{1.5em} ...}
%
% That is enough for two digits and not for three. At 10.95pt bold, "100" is
% 2.46pt wider than the box, so the number runs into its own title and TeX
% reports an overfull \hbox for every such line. A track notices nothing until
% it crosses a hundred chapters -- and then the bug arrives all at once, one
% line per chapter, which is exactly how it surfaced here.
%
% Below is `book.cls`'s own \l@chapter with a single number changed: 1.5em
% becomes 2.8em, which holds four bold digits with room to spare. A course that
% reached 100 chapters can reach 1000, and the fix should not need revisiting
% then. The section and subsection indents move by the same amount so the
% columns still line up under the widened chapter number.
\makeatletter
\renewcommand*\l@chapter[2]{%
  \ifnum \c@tocdepth >\m@ne
    \addpenalty{-\@highpenalty}%
    \vskip 1.0em \@plus\p@
    \setlength\@tempdima{2.8em}%
    \begingroup
      \parindent \z@ \rightskip \@pnumwidth
      \parfillskip -\@pnumwidth
      \leavevmode \bfseries
      \advance\leftskip\@tempdima
      \hskip -\leftskip
      #1\nobreak\hfil
      \nobreak\hb@xt@\@pnumwidth{\hss #2%
                                 \kern-\p@\kern\p@}\par
      \penalty\@highpenalty
    \endgroup
  \fi}
\renewcommand*\l@section{\@dottedtocline{1}{2.8em}{3.2em}}
\renewcommand*\l@subsection{\@dottedtocline{2}{6.0em}{4.1em}}
\makeatother

% ---------------------------------------------------------------------------
% Line-breaking tolerance (HL-C109).
%
% These books mix ordinary prose with tokens TeX cannot break: lesson ids
% (ES-C14-practice), slash-separated form lists (fui/fuiste/fue), and Latin
% etymons carrying macrons that a modern language's hyphenation patterns will
% not touch. Inside a narrow measure -- a quote block or a list item -- one such
% token forces the previous line to stretch, and TeX reports an underfull hbox.
%
% \tolerance and \emergencystretch change how a paragraph is BROKEN, not how it
% is reported: TeX may accept a looser line early enough to find a break
% sequence that avoids a very loose one, and gets a final pass with extra
% stretchability instead of an overfull line. \hbadness and \hfuzz are
% deliberately left at their defaults -- raising those would silence the
% warnings rather than fix the paragraphs, and the CI log scanner would then be
% reading a muted log.
\tolerance=2000

% Callout boxes are set ragged right (HL-C88).
%
% Every callout runs at a narrower measure than the body, and its content is the
% densest part of the book: bold Spanish forms, Latin etymons carrying macrons,
% and arrows between them. None of those hyphenate under English patterns, so a
% justified line inside a callout has almost nowhere to break and must stretch
% the few interword spaces it has -- which is precisely an underfull hbox.
%
% Ragged right removes the stretch rather than hiding it: the line simply ends
% where the last word ended. Block callouts set ragged right is ordinary
% typography, and it is what \chapteropening already does above.
\tcbset{before upper=\raggedright}
\emergencystretch=2em

% ---------------------------------------------------------------------------
% Running heads (HL-C109).
%
% Chapter titles here are deliberately descriptive -- "Answering, Meeting, and
% Three Ways to Play" -- and `book.cls` puts the whole title into the verso
% running head as one unbreakable line. Past a certain length it simply
% overflows the text block, once per verso page for the length of the chapter,
% which is how a single long title produced four identical overfull boxes.
%
% The head shows "Chapter N" instead. That keeps it useful and stable, and it
% cannot overflow no matter what a chapter is called. Spanish has done this for
% a while; this is the same rule, applied to every track.
\renewcommand{\chaptermark}[1]{\markboth{\chaptername\ \thechapter}{}}
\renewcommand{\sectionmark}[1]{\markright{\thesection}}
